% Source: https://artofproblemsolving.com/community/c3h1459736_prmo_problem
% Stable identity: aops-topic-1459736-q01
We need $\dfrac{K(K+1)}{2}=N^2$. The numbers $K$ and $K+1$ are coprime and one of them is even.

\textbf{Case 1: $K$ even.} Then $\dfrac {K}{2}$ and $K+1$ are coprime with product $N^2$, so both are squares: $K=2a^2$, $K+1=b^2$, $N=ab$, with $b^2=2a^2+1$. Since $b>\sqrt2\,a$, we get $100>N=ab>\sqrt2\,a^2$, so $a\le8$. For $a=1,\dots,8$ the values $2a^2+1$ are $3,9,19,33,51,73,99,129$; only $9$ is a square ($a=2$, $b=3$). This gives $K=8$, $N=6$.

\textbf{Case 2: $K$ odd.} Then $K$ and $\dfrac{K+1}{2}$ are coprime squares: $K=a^2$, $K+1=2b^2$, $N=ab$, with $a^2=2b^2-1$. Here $a\ge b$, so $N\ge b^2$ and $b\le9$. For $b=1,\dots,9$ the values $2b^2-1$ are $1,7,17,31,49,71,97,127,161$; the squares are $1$ ($b=1$) and $49$ ($b=5$, $a=7$). These give $K=1$, $N=1$ and $K=49$, $N=35$.

So $K\in\boxed{\{1,\,8,\,49\}}$ (with $1=1^2$, $36=6^2$, $1225=35^2$). The next solution, $K=288$, has $N=204>100$.
